% ======================================================================
% Ontology of Continua -- Core 1.3.3
% Cross-K module: crossk_k6_k7.tex
% Cross-level relations between K6 and K7
% ======================================================================

\subsubsection{Межуровневая структура K6–K7}
\label{sec:crossk-k6-k7}

Переход $K_6 \rightarrow K_7$ обозначает возникновение социального
континуума из когнитивного. Если $K_6$ поддерживает концептуальные структуры,
смысл и внутренние модели, то $K_7$ вводит:
\begin{itemize}
  \item социальные роли и нормы как новые оси,
  \item потоки коммуникации и общих символических структур,
  \item коллективные пороги координации и ожиданий,
  \item институциональные циклы и стабилизирующие обратные связи,
  \item паттерны кооперации, конфликта и нормоприменения,
  \item общую область $\Omega(K_7)$ социально допустимых конфигураций.
\end{itemize}

Этот переход соответствует возникновению коллективной организации,
не сводимой к индивидуальной когниции.

% ----------------------------------------------------------------------
\subsubsection{1. Задействованные уровни и их роли}

\paragraph{K6 (когнитивный континуум).}
$K_6$ предоставляет:
\begin{itemize}
  \item концептуальные оси $A^{c}$,
  \item когнитивные потенциалы $P^{c}$,
  \item пороги $\Theta^{c}$ для когерентности и противоречивости,
  \item потоки внимания, памяти, интерпретации $J^{c}$,
  \item циклы понимания и обучения,
  \item меру когнитивной когерентности $k_6$.
\end{itemize}

Однако у $K_6$ отсутствуют:
\begin{itemize}
  \item правила координации между агентами,
  \item общие символические нормы,
  \item популяционные пороги устойчивости,
  \item институциональная инерция.
\end{itemize}

\paragraph{K7 (социальный континуум).}
$K_7$ вводит:
\begin{itemize}
  \item оси $A_{\mathrm{soc}}$ (роли, нормы, статус, идентичность),
  \item потенциалы $P_{\mathrm{soc}}$ (доверие, авторитет, легитимность),
  \item социальные пороги $\Theta_{\mathrm{soc}}$ (координация, сплоченность,
        принуждение, коллапс),
  \item потоки $J_{\mathrm{comm}}$ (коммуникация, сигнализация, подражание),
  \item циклы $C_{\mathrm{inst}}$ (правило–действие–последствие–исправление),
  \item меру $k_7$, описывающую социальную стабильность.
\end{itemize}

$K_7$ — первый континуум, в котором ожидания и нормы регуляторно влияют на
поведение множества агентов.

% ----------------------------------------------------------------------
\subsubsection{2. Совместные и наследуемые оси и пороги}

\paragraph{Наследование.}
Концептуальные различия из $K_6$ служат подложкой для общих символов:
\[
A^{c}(K_6) \rightarrow A_{\mathrm{soc}}(K_7).
\]
Внутренние модели превращаются в основу общих нарративов.

\paragraph{Новые оси.}
$K_7$ вводит:
\begin{itemize}
  \item $A_{\mathrm{role}}$ — поведенческие ожидания,
  \item $A_{\mathrm{norm}}$ — допустимые/недопустимые действия,
  \item $A_{\mathrm{status}}$ — социальная дифференциация,
  \item $A_{\mathrm{collective}}$ — структуры групповой идентичности.
\end{itemize}

Эти оси отсутствуют в $K_6$.

\paragraph{Пороги.}
$\Theta_{\mathrm{soc}}$ включает:
\begin{itemize}
  \item пороги сплоченности (минимальное доверие или плотность коммуникаций),
  \item пороги противоречия (несовместимость норм),
  \item пороги принуждения (устойчивость санкций),
  \item пороги коллапса (отказ институциональных циклов).
\end{itemize}

Нарушение $\Theta^{c}$ разрушает когерентность на когнитивном уровне и
предотвращает возникновение $K_7$.

% ----------------------------------------------------------------------
\subsubsection{3. Межуровневые потоки, циклы и напряжения}

\paragraph{Потоки.}
Социальные потоки $J_{\mathrm{comm}}$ расширяют когнитивные потоки:
\[
J_{\mathrm{comm}} = (J_{\mathrm{signal}},\ J_{\mathrm{language}},\
                     J_{\mathrm{imitation}},\ J_{\mathrm{coord}}).
\]
Они перераспределяют смысл между агентами и стабилизируют коллективные состояния.

\paragraph{Циклы.}
$K_7$ характеризуется институциональными циклами:
\[
C_{\mathrm{inst}}:\ \text{rule} \rightarrow \text{action}
\rightarrow \text{consequence} \rightarrow \text{correction}
\rightarrow \text{rule}.
\]
Общество существует, только если остаётся хотя бы один такой стабилизирующий цикл.

\paragraph{Напряжение.}
Межуровневое напряжение $T_{6,7}$ измеряет несовместимость между личными
концептуальными моделями и коллективными структурами:
\[
T_{6,7} = f(A_{\mathrm{soc}} - A^{c},\ P_{\mathrm{soc}} - P^{c},\
          J_{\mathrm{comm}} - J^{c},\ \Theta_{\mathrm{soc}} - \Theta^{c}).
\]
Если $T_{6,7}$ превышает размерностный порог, коллективная организация
становится невозможной и система схлопывается к индивидуальной когниции.

% ----------------------------------------------------------------------
\subsubsection{4. Условия рождения и исчезновения между уровнями}

\paragraph{Рождение \texorpdfstring{$K_7$}{K_7}.}
$K_7$ возникает при:
\[
T_6 > \Theta_{6,\mathrm{dim}}
\quad\text{и}\quad
A_{\mathrm{soc}} \in M_7 \setminus A(K_6).
\]
Это соответствует формированию устойчивых социальных норм и ролей.

Расширение пространства состояний:
\[
\Omega(K_7) = \Omega(K_6) \cup \Delta\Omega_{\mathrm{soc}}.
\]

\paragraph{Распространение гибели.}
Если $\Theta^{c}$ нарушена (когнитивный коллапс), $K_7$ становится невозможной.

Если $\Theta_{\mathrm{soc}}$ нарушена (потеря сплоченности, разрушение циклов),
система схлопывается до индивидуальной когниции:
\[
\Omega(K_7) \rightarrow \Omega(K_6).
\]

$K_6$ выживает, если не нарушены его собственные пороги.

% ----------------------------------------------------------------------
\subsubsection{5. Концептуальные примеры}

\paragraph{Возникновение норм.}
Повторяющиеся когнитивные интерпретации по всем агентам стабилизируются в
коллективные нормы, образуя новые оси в $A_{\mathrm{soc}}$.

\paragraph{Институциональная устойчивость.}
Когда институциональный цикл $C_{\mathrm{inst}}$ замыкается с минимальной
дисперсией:
\[
\oint_C dA_{\mathrm{soc}} \approx 0,
\]
формируется стабильная социальная подсистема.

\paragraph{Сбой.}
Если плотность коммуникации падает ниже порога сплоченности:
\[
J_{\mathrm{comm}} < \Theta_{\mathrm{soc,coh}},
\]
коллективная структура распадается.

Следовательно, переход K6→K7 можно сжато описать как возникновение
общих норм, коммуникационных потоков и институциональных циклов из
индивидуальной когниции.
